\chapter{Implementazione della pipeline RAG}

In questo capitolo tratteremo estensivamente la fase implementativa del progetto e, come preannunciato, quali siano state le principali sfide da risolvere durante lo sviluppo di quest'ultimo.

\section{Processamento preliminare della query dell'utente e disambiguazione}
Si è pensato fin da subito di processare la richiesta dell'utente tramite un LLM (detto "\textit{guardrail}"), richiedendo a quest'ultimo, mediante un prompt adeguato, di classificare la query come \texttt{ACCEPTED} o \texttt{REJECTED}. Nel caso in cui dovesse essere rifiutata, l'utente viene notificato con il messaggio "\texttt{Your query appears to be meaningless or invalid. Please ask a clear, answerable question}". 

In seguito è stata introdotta una terza possibile categorizzazione della query dell'utente, ossia \texttt{AMBIGUOUS}, nel caso in cui il soggetto della richiesta dovesse risultare omesso oppure dovessero essere individuate entità con
molteplici significati (\textit{i.e.} la parola "mercurio" può intendere l'elemento chimico, ma anche il pianeta). In tal caso il sistema sospende il retrieval e richiede delucidazioni all’utente tramite
una query opportuna generata dinamicamente.

Nelle ultime versioni del progetto si è affidato all'LLM \textit{guardrail} l'ulteriore compito di identificare e scartare in maniera sistematica richieste malevole da parte dell'utente, come, ad esempio, tentativi di prompt injection, \textit{persona replacement}, e/o manipolazione emotiva.

In realtà ogni LLM della pipeline è istruito in tal senso, in modo tale da non far ricadere \textit{in primis} la responsabilità interamente sul \textit{guardrail} e in secondo luogo di permettere esclusivamente alle query legittime, o più in generale al sottoinsieme delle richieste accettabili di queste ultime, di essere processate dal sistema e di ricevere una risposta.

\section{Assegnazione della query a un dominio e fallback}
Una sfida significativa consiste nel \textit{routing} di una richiesta ricevuta. In particolare, per ogni richiesta in input alla pipeline, è necessario porsi la seguente domanda: "qual è il dominio più adatto a soddisfare la query dell'utente?".

Non è semplice trovare un procedimento deterministico che risponda in maniera precisa a questa domanda. Pertanto, anche in questo caso ci si è affidati a un LLM (difatti il \textit{guardrail} stesso) per indirizzare la query dell'utente verso il dominio ritenuto più opportuno tra i seguenti: \href{it.wikipedia.org}{it.wikipedia.org}; \href{www.raiplaysound.it}{www.raiplaysound.it}; \href{www.ipsos.com}{www.ipsos.com}; \href{www.marvel.com}{www.marvel.com}.

Si noti che in alcuni casi potrebbe accadere che la richiesta effettuata dall'utente non possa trovare risposta in alcuno dei domini sopra riportati. Per questo motivo, in un secondo momento, si è permesso all'LLM di decidere di assegnare un dominio generico (individuato tramite "⁎") alla query ricevuta come meccanismo di fallback.

Quando ciò accade si tenta di cercare le informazioni necessarie per rispondere alla domanda su un dominio qualunque, effettuando un parsing generico di pagine web ritenute pertinenti.

A un secondo LLM viene dunque affidato il compito di sintetizzare la domanda dell'utente in poche \textit{keyword} che costituiranno la query per motore di ricerca.

La peculiarità di questo procedimento consiste nel fatto che l'output dell'LLM dipenda anche dalla {chat history}, vale a dire la cronologia delle interazioni precedenti con l'utente

In altre parole, se è stato richiesto all'utente un chiarimento sulla query iniziale, sia essa ad esempio "parlami di Mercurio" e l'utente ha disambiguato con "mi riferisco al pianeta", l'LLM non proporrà come \texttt{search\_query} "pianeta", ma "pianeta mercurio". Il sistema è pertanto in grado di mantenere il contesto nonché filo logico della conversazione, anche qualora l'utente dovesse disconnettersi, grazie alla memorizzazione persistente sul database.

A partire dalle ultime implementazioni del progetto, si è deciso di fare in modo che l'LLM, data la query in input, non ritornasse esclusivamente la singola query per il search engine, ma una coppia <\texttt{search\_query}, \texttt{user\_query}>, in modo tale da poter passare la \texttt{user\_query} all'LLM responsabile di fornire la risposta all'utente, senza perdere informazioni sulla domanda fatta originariamente. 

\section{Ricerca degli URL e parsing delle pagine web}
Data una query dell'utente e il dominio assegnato a quest'ultima è necessario trovare gli URL delle pagine che contengono le informazioni necessarie a rispondere ad essa.

Ci si è ben presto resi conto che ottenere in maniera programmatica i risultati di una ricerca su web da un search engine non fosse semplice. In effetti, tutte le API rese disponibili per questo scopo si sono rivelate essere a pagamento oppoure con limiti di utilizzo particolarmente stringenti.

\subsection{DDGS - Duck Duck Go Search}
Una prima soluzione alla quale si è giunti consistenva nell'utilizzo della libreria \texttt{DDGS}, basata sul motore di ricerca \textit{Duck Duck Go}. Questa soluzione si è rivelata tutto fuorché ottimale poiché dopo poche ricerche il motore ha segnalato le nostre richieste come sospette e ha cominciato a non restituire più alcun risultato.

\subsection{SearXNG e qwant}
Un tentativo successivo ha visto l'impiego di \texttt{SearXNG}, un \textit{metasearch engine} open-source che aggrega risultati di ricerche web dai principali servizi di ricerca disponibili (\textit{Google}, \textit{Bing}, \textit{Brave}, etc.). L'unico servizio di ricerca che si è rivelato essere realmente utilizzabile è stato \texttt{qwant}, dal momento che tutti gli altri motori hanno cominciato a bloccare immediatamente le richieste provenienti dal nostro IP.

Purtroppo, anche in questo caso, dopo alcuni giorni di testing, l'unico motore di ricerca a nostra disposizione ha sospeso il nostro IP lasciandoci senza alcuna alternativa per la ricerca di pagine web da cui estrarre i dati per la nostra pipeline.

\subsection{Custom web scraper per Startpage.com}
In seguito a un'ampia riflessione si è giunti a una nuova soluzione: parsare le pagine stesse contenenti i risultati di una determinata query effettuata su un search engine. Trovare un modo concreto di farlo si è rivelato essere non banale.

Infatti, l'analisi del DOM di una pagina restituita da una ricerca su \textit{Google} ha ben presto evidenziato una problematica sostanziale: ogni singolo \textit{div} o \textit{span} contenente i risultati della ricerca era caratterizzato da classi e \textit{id} generati in modo del tutto casuale per ogni nuova ricerca, rendendo impossibile estrarre i dati in modo automatizzato.

Per puro caso, tuttavia, siamo stati in grado di trovare un motore di ricerca privato che fa utilizzo di una ben nota API di ricerca a pagamento: \texttt{SerpAPI}. Tale motore di ricerca è \href{www.startpage.com}{Startpage} e, al contrario di \textit{Google}, non implementa meccanismi \textit{anti-crawling} nelle pagine restituite in seguito a una query.

Pertanto, usufruendo ancora una volta della libreria \texttt{Crawl4AI} e modificando leggermente il progetto dell'esonero (\texttt{Minerva Web Parser}) in modo tale da introdurre un nuovo parser (\texttt{StartpageSearchEngineParser}) capace di restituire, data una query, un dominio di interesse e un intero \textit{k}, i primi \textit{k} risultati di tale ricerca (senza duplicati) su \texttt{Startpage}. Inoltre, tramite l'aggiunta di un semplice endpoint \texttt{/get\_query\_results} abbiamo difatti creato una nostra API personalizzata al fine di permettere lo scraping e il \textit{retrieval} degli URL risultanti da una ricerca su web.

Si potrebbe argomentare che i risultati restituiti possano non essere affidabili quanto quelli forniti da \textit{Google}, ma la convenienza della soluzione appena presentata consta proprio del fatto che \texttt{SerpAPI} stesso ottenga i propri risultati a partire da quest'ultimo.

\subsection{Estrazione dei dati dagli URL selezionati}
A questo punto, dati gli URL ottenuti da una query al motore di ricerca e il dominio target è immediato estrarne in contenuti in formato markdown.

Tuttavia, bisogna tenere presente che a una determinata query \textit{q}, in caso di fallback, può essere assegnato un dominio generico "⁎". Anche in questo caso abbiamo dovuto modificare leggermente il progetto di esonero, aggiungendo un secondo parser, \texttt{GenericParser}, adatto ad estrarre in modo approssimativo i contenuti di una qualunque pagina web.

Naturalmente, risulta impossibile effettuare quest'operazione senza introdurre una quantità minima di rumore (dati inutili) dovuto alla sostanziale differenza che si può avere tra pagine di domini diversi. Si tratta di un \textit{trade-off} che abbiamo ritenuto accettabile.

\subsection{Supplemento dei dati ottenuti tramite generic parsing}
Per permettere all'LLM che risponde alla query dell'utente di disporre della massima quantità di informazioni relative all'argomento della conversazione, ogni ricerca per il dominio \textit{X} viene supplementata con i primi \textit{k} risultati della stessa ricerca effettuata su domini generici, purché il dominio scelto dal guardrail non sia già "⁎".

Questo equivale ad effettuare due query al \textit{search engine}, dove la prima è "\texttt{site:X \{query\}}" mentre la seconda consiste semplicemente nel ricercare "\texttt{\{query\}}" senza utilizzare l'operatore "\texttt{site:}". 

\section{Selezione dei contenuti rilevanti, generazione e gestione del contesto per l'LLM}
In questa sezione verrà trattato il processo tramite cui si è in grado di scegliere quali testi (\texttt{parsed\_text}) e quali porzioni di questi ultimi sono da considerare rilevanti e dunque da fornire all'LLM che risponde alla query dell'utente.

Si tratta fondamentalmente della sfida più ardua nella costruzione della pipeline RAG: il contesto dell'LLM non è infinito e va dunque gestito opportunamente in modo tale da non saturarlo.

Difatti, se decidessimo di fornire pagine intere come fonti di riferimento al modello, non solo rischieremmo di esaurire il \textit{context window} ma introdurremmo un notevole overhead in fase di risposta, diminuendo notevolmente il throughput della pipeline.

L'idea alla base è quella di stilare una classifica dei pezzi (o "\textit{chunk}") di testo che contengono, con maggiore probabilità, la risposta alla domanda dell'utente.

Una volta ottenuta la classifica (o "ranking") di questi ultimi, è sufficiente prendere i primi \textit{k} chunk, i quali andranno a costituire il \texttt{query\_context} dell'LLM, ossia l'insieme di testi di riferimento nonché informazioni che il modello utilizzerà per rispondere all'utente, purché in essi vi sia la risposta alla query.

In caso contrario, il modello risponderà semplicemente con "\texttt{Mi dispiace, ma non ho abbastanza informazioni per rispondere a questa domanda}".

Di seguito è proposto un resoconto degli approcci utilizzati, in fase di sviluppo del progetto, per raggiungere l'obiettivo appena descritto. 
\subsection{Chunking manuale e token count}
Un primo approccio \textit{naive} è stato quello di dividere manualmente il \texttt{parsed\_text} ottenuto da ciascun URL in \textit{chunk} di lunghezza arbitraria.

In questa prima soluzione proposta, i \textit{chunk} venivano classificati per ordine di count decrescente di \texttt{query\_token} contenuti in essi.

Purtroppo, questa soluzione semplice ed estremamente conveniente comportava in realtà risultati piuttosto scarsi.

\subsection{Chunking manuale e bi-encoding, ranking mediante cosine similarity}
Un altro difetto del primo approccio consisteva nel fatto di tagliare arbitrariamente il chunk una volta raggiunti \textit{k} caratteri. Questo poteva comportare il troncamento di paragrafi, frasi o addirittura parole.

In questa seconda iterazione di selezione abbiamo voluto rifinire il processo di chunking, tenendo conto di queste problematiche, ad esempio, aggiungendo il testo contenuto tra una newline e la successiva, fino al raggiungimento della lunghezza massima per un chunk.

In realtà anche questo modo di procedere non era perfetto: due frasi separate da newline e appartenenti allo stesso paragrafo potevano essere divise in chunk diversi, comportando la dispersione di informazioni.

A questo punto sia query che chunk estratti venivano convertiti in vettori \textit{n}-dimensionali, dove \textit{n} dipende dal modello di embedding utilizzato.

Il modello utilizzato in questa fase è stato inizialmente \texttt{nomic-embed-text}. Ogni porzione di testo da tradurre veniva tradotto in un vettore mediante l'opportuna chiamata API.

Infine, per ciascuna coppia <\texttt{query\_vec}, \texttt{chunk\_vec}> veniva calcolata la metrica di cosine similarity, utilizzata successivamente come score per i singoli \textit{chunk}. La lista di \textit{chunk} veniva riordinata sulla base di questi ultimi e tutti i \textit{chunk} con punteggi inferiori a una certa soglia, ossia \texttt{COSINE\_SIMILARITY\_THRESHOLD}, venivano automaticamente scartati.

Ora, si selezionavano i primi \textit{k} nella lista e si costruiva il contesto per l'LLM, fino al raggiungimento di una lunghezza massima predefinita, cioè \texttt{MAX\_OUTPUT\_LENGTH}.

\subsection{RecursiveCharacterTextSplitter, bi-encoding e cross-encoding}
Nell'implementazione finale di questa componente della pipeline si è deciso di sostituire il chunking manuale dei testi con una funzione utility, \texttt{RecursiveCharacterTextSplitter} della libreria \texttt{LangChain}.

Tale modifica permette di suddividere in modo preciso il testo in \textit{chunk}, senza perdere informazioni e mantenendo una soglia limite di \textit{k} caratteri.

Ancora una volta si procede al bi-encoding di query e chunk, tuttavia con un modello più leggero (\texttt{intfloat/multilingual-e5-small}, fornito dalla libreria \texttt{TextEmbedder}) e in esecuzione direttamente nel container del \textit{backend} (nell'iterazione precedente il modello si trovava in esecuzione su un container separato, introducendo maggiore overhead).

In aggiunta a ciò, \texttt{TextEmbedder} permette l'embedding di testo in \textit{batch}, operazione indubbiamente più efficiente dell'embedding iterativo di ogni singolo \textit{chunk} ottenuto dai testi estratti nella fase precedente.

Il passaggio successivo rimane invariato, difatti, una volta effettuato l'encoding si stila una lista sulla base dello score in termini di \textit{cosine similarity} e si scartano tutti i chunk al di sotto della soglia stabilita.

Tuttavia, un alto score di \textit{cosine similarity} garantisce solo la vicinanza semantica del \textit{chunk} per cui è stato calcolato alla query dell'utente, ma non assicura che quest'ultimo contenga veramente la risposta alla domanda posta.

Per questo motivo, si è deciso di implementare un \textit{re-ranking} dei primi \textit{k} chunk così ottenuti per costruire la "graduatoria" finale. Si utilizza pertanto un \textit{cross-encoder} (nel nostro caso \texttt{ms-marco-MiniLM-L-6-v2}), in grado di stabilire quali siano realmente i \textit{chunk} contenenti le informazioni utili a rispondere alla \texttt{user\_query}.

Dopo l'applicazione, anche in questo caso, di una soglia di punteggio minimo, ossia \texttt{CROSS\_ENCODING\_SIMILARITY\_THRESHOLD},  si aggiungono i \textit{chunk} al \texttt{query\_context}, fino al raggiungimento della lunghezza massima per quest'ultimo.

Tale implementazione assicura che, qualora le informazioni estratte, o quantomeno una porzione di esse, dovesse contenere la risposta alla domanda dell'utente, quest'ultima sarà certamente inserita nel contesto dell'LLM, permettendo dunque a quest'ultimo di fornire la risposta corretta all'utente.

\section{Iniezione del \texttt{query\_context} nel contesto dell'LLM e generazione della risposta finale per l'utente}
Si tratta difatti del passaggio più semplice tra tutti quelli visti finora. L'unica difficoltà consiste nell'effettuare "\textit{grounding}" dell'LLM, vale a dire evitare, attraverso \textit{prompt engineering}, che il modello utilizzi informazioni di cui è già in possesso per rispondere alla domanda dell'utente.

L'obiettivo è proprio avere un modello che sia in grado di rispondere alla \texttt{user\_query} se e solo se la risposta è effettivamente contenuta nelle fonti ad esso fornite.

Ora è sufficiente iniettare la \texttt{user\_query} nel prompt dell'LLM assieme al \texttt{query\_context} e attendere la sua risposta.

\subsection{Stima dell'affidabilità ed assegnazione da parte dell'LLM del relativo score}
Al termine della propria risposta, una regola ferrea stabilita in fase di prompting impone all'LLM di emettere obbligatoriamente una breve sezione denominata "affidabilità", costituita da un breve commento sulla propria risposta, in cui il modello specifica quante delle informazioni fornite derivino direttamente dal testo e quanto queste ultime siano da ritenere pertanto attendibili.

Il modello fornisce, accanto al commento, un punteggio indicativo (da \textit{1} a \textit{5}) dell'indice di affidabilità (o "\textit{reliability}") della risposta data all'utente.

Questo garantisce \textit{in primis} trasparenza e permette inoltre all'utente di stabilire in modo indicativo in che porzione attenersi alla risposta appena ricevuta, sebbene sia sempre opportuno che quest'ultimo verifichi indipendentemente tali informazioni.

\subsection{Trasparenza ed enumerazione delle fonti di riferimento}
A tal riguardo, in fondo alla risposta del modello sono sempre indicate le fonti utilizzate da quest'ultimo per rispondere alla domanda dell'utente.

Questo permette permette, tra le altre cose, una verifica celere di quanto riportato nella risposta generata dall'LLM.